\subsection{Hypothesis and Test Procedures}
\hrulefill

\subsubsection*{Hypothesis Testing}
\paragraph*{Definition}
A statistical hypothesis is a claim about a population parameter. The null hypothesis, denoted by $H_0$, is the claim that is initially
assumed to be true. The alternative hypothesis, denoted by $H_a$, is an assertion that is contradictory to $H_0$. Null hypothesis is 
the status quo, aka the default assumption is absence of compelling evidence. The alternative hypothesis is what we want to find evidence for.

At the conclusion of a hypothesis test, we either reject $H_0$ in favor of $H_a$ or fail to reject $H_0$. We never accept $H_0$ as true, 
we only fail to reject it. An analogue to this is a court trial, where the defendant is assumed innocent until proven guilty. 
If there is enough evidence to convict, the defendant is found guilty (reject $H_0$). If there is not enough evidence, the defendant is
found not guilty (fail to reject $H_0$).

\paragraph*{Errors}
There are two possible errors that can be made in hypothesis testing:
\begin{itemize}
    \item Type I Error: Rejecting the null hypothesis when it is actually true. (Finding a guilty defendant innocent)
    \item Type II Error: Failing to reject the null hypothesis when it is actually false. (Finding an innocent defendant guilty)
\end{itemize}

\begin{itemize}
    \item We want to keep these two types of error rates as small as possible
    \item However, it is impossible to minimize both simultaneously
    \item The choice of rejection region determines the probabilities of both a type I and type II error
\end{itemize}

\paragraph*{Significance Level}
The probability of a type I error is called the significance level of the test, denoted by $\alpha$. The probability of a type II error 
is usually denoted by $\beta$. We agree on a maximum acceptable value for $\alpha$ before conducting the test, commonly 0.05, 0.01, or 0.005.
The quantity $1 - \beta$ is called the power of the test, which is the probability of correctly rejecting a false null hypothesis. Good tests
have high power.

\paragraph*{Components of a Statistical Test}
\begin{itemize}
    \item A \textbf{test statistic} (based on a random sample from the population under question) that summarized the evidence against the null hypothesis
    \item A \textbf{p-value} is a probability that measures how likely the value of the test statistic is, assuming that $H_0$ is true
    \item A \textbf{rejection region} tells us for which values of the test statistic we reject $H_0$
    \item A \textbf{Conclusion} is a decision of the form either Reject $H_0$ or Fail to Reject $H_0$
    \item A \textbf{significance level} $\alpha$ which represents the maximum allowable risk of rejection $H_0$ in favor of $H_a$ when $H_0$ is true.
\end{itemize}

\paragraph*{p-values}
One way to report the results of a hypothesis test is to determine if the test statistic value falls into the rejection region and subsequently, whether or not the 
null hypothesis is rejected at a specified level of significance. This yes/no decision does not convey any information about how soundly the 
null hypothesis was rejected. Where in the rejection region did the observed test statistic fall? Note that the p-value is calculated assuming
that the null hypothesis is true.

\paragraph*{Definition}
Suppose the null hypothesis $H_0$ is true. The p-value is the probability of observing a test statistic value 
at least as contradictory to $H_0$ as the computed value due to sampling variability. In other words, the p-value is the probability of obtaining a test statistic
value as extreme or more extreme than the one observed, assuming that $H_0$ is true.

\subsubsection*{Rejection Regions}
\paragraph*{Definition:} If a p-value is smaller than the significance level $\alpha$, then the 
corresponding value of the test statistic falls into the rejection region and the null hypothesis is rejected in favor of the alternative hypothesis.
\[
    p\text{-value} \leq \alpha \rightarrow \text{reject } H_0
\]
If the p-value is large, then it is quite likely to see data as extreme as that observed, assuming that $H_0$ is true. Thus, there is not enough evidence to reject $H_0$.
\[
    p\text{-value} > \alpha \rightarrow \text{fail to reject } H_0
\]

\paragraph*{Example:}
A kid is supposed to bush their teeth every night. One night, the parent asks the kid if they have already brushed their teeth.
The kid says yes. The parent wants to test the hypothesis that the kid has brushed their teeth. The null hypothesis is that the kid has brushed their teeth,
$H_0: \text{brushed teeth}$. The alternative hypothesis is that the kid has not brushed their teeth, $H_a: \text{not brushed teeth}$.
The test statistic is whether or not the toothbrush is wet. If the toothbrush is wet, then the p-value is high, and the parent fails to reject $H_0$.
If the toothbrush is dry, then the p-value is low, and the parent rejects $H_0$ in favor of $H_a$.

\paragraph*{Example:}
Suppose $H_0 : \mu = 5$ with alternative $H_a : \mu \neq 5$. (a) Find the rejection region for the test statistiv for $\alpha = 0.05$.
(b) Compute the p-value for an observed test statistic value of 2.14.

\subsubsection*{Tests for a Population Mean}
\paragraph*{Case 1:} Sampling from a Normal Population with known $\sigma$. Let $X_1, X_2, \ldots, X_n$ be a random sample from a normal population with
unknown mean and known variance $\sigma^2$. According to the Central Limit Theorem, 
the sample mean $\bar{X}$ has a Normal distribution with mean $\mu$ and variance $\frac{\sigma^2}{n}$.
\[
    Z = \frac{\bar{X}-\mu_0}{\frac{sigma}{\sqrt{n}}} \sim \mathcal{N}(0,1)
\]
Suppose we want to test the null hypothesis that $H_0 : \mu = \mu_0$ (some specified value). In this case our test statistic is:
\[
    Z = \frac{\bar{X}-\mu_0}{\frac{\sigma}{\sqrt{n}}}
\]
The distribution of the test statistic is known if $H_0$ is true. This allows us to find rejection
regions for specified levels of $\alpha$ or to compute p-values for certain values $z$ of 
the test statistic.

\paragraph*{Case 2:} 
If a sample is large $(n \geq 40)$, then the population standard deviation $\sigma$ may be replaced by the point esitmate $s$ 
without changing the distribution of the sample mean, meaning our test statistic becomes:
\[
    Z = \frac{\bar{X}-\mu_0}{\frac{s}{\sqrt{n}}}
\]
This is an approximation that is valid for large samples due to the Central Limit Theorem.

\paragraph*{Example:} 
The daily yield for a chemical plant has avereaged 880 tons for several years. The quality control manager
wants to know if this average has changed. She randomly selects 50 days and records and average yield of
871 tons with a standard deviation of 21 tons. Perform a test with sigificance level $\alpha = 0.05$.
\begin{align*}
    t = \frac{871 - 880}{\frac{21}{\sqrt{50}}} = -3.03\\
    \text{Degrees of Freedom} = n - 1 = 49\\
    t' = \pm 2.009\\
    t < -t' \rightarrow \text{Reject } H_0
\end{align*}

\subsubsection*{Test for a Population Proportion}
Let $p$ denote the proportion of individuals in a popoulation who possess a certain characeristic.
A random sample of size $n$ is selected from the population and the proportion $\hat{p}$ of successes
in the sample is observed. We want to use this quantity to decide if we will reject a statement about
population parameter $p$.

\paragraph*{Large Sample Test for a Population Proportion}
If the sample size $n$ is large $(np_0 \geq 10$ and $n(1-p_0) \geq 10)$, then the test statistic is:
\[
    H_0 : p = p_0
\]
has approximately a Normal distribution. Test Statistic:
\[
    Z = \frac{\hat{p} - p_0}{\sqrt{\frac{p_0(1 - p_0)}{n}}}
\]
where $p_0$ is the hypothesized value of the population proportion under the null hypothesis.

\paragraph*{Example:}
Consider Prevnar, a vaccine for meningitis usually given to infants. In a clinical trial, the vaccine was given to 710 children,
of whom 72 experiences a loss of appetitie. Competing medications cause about 13.5\% of 
children to experience a loss of appetitie. Can we conclude that the percentage of children experiencing a loss of 
appetite from Prevnar is significantly less than for other medications? Use a significance level of $\alpha = 0.01$.
\begin{align*}
    H_0 : p = 0.135\\
    H_a : p < 0.135\\
    \hat{p} = \frac{72}{710} = 0.1014\\
    Z = \frac{0.1014 - 0.135}{\sqrt{\frac{0.135(1-0.135)}{710}}} = -2.36\\
    Z' = -2.33 \text{ ~1000 DOF } \alpha = 0.01\\
    Z < Z' \rightarrow \text{Reject } H_0
\end{align*}